\section{Theorem G: the necessary frozen gap area}\label{sec:gap-area}

The paid inequality of \cite[Theorem~3.5]{SourceCovers2026} accepts any credit that is
at most the frozen point area at each legal phase and column. Here the credit is the
frozen point area itself, and the limits are taken as for the reciprocal-cover bound
\cite[Theorem~4.6]{SourceCovers2026}.

\begin{definition}[Frozen gap-area means]\label{def:gap-area-means}
\formal{GapAreaMemo/GapArea.lean}{SourceCovers.phaseGapArea,SourceCovers.phaseGapArea_bounds,SourceCovers.meanGapArea,SourceCovers.AdmissibleGeometry.prefixMeanGapArea,SourceCovers.NearFullSequence.JointProfile.prefixGapArea,SourceCovers.NearFullSequence.JointProfile.gapAreaLimit,SourceCovers.NearFullSequence.JointProfile.gapAreaLimit_bounds}
Fix a geometry, a $\mathrm Q_2$ selected core, an exceptional set, $M\ge2$ and a
protected set of at most $R<M$ selected families. For a residue $\rho<M$ and a column
$j<d$, the frozen point area $\gaparea(S_\rho(s_j))$ lies in $[0,1/2]$. Let $\meanarea$
be its arithmetic mean over the legal phases $\rho$ and all columns $j$. For an admissible geometry and an order of its original core by
nonincreasing hop count, $\meanarea_{J,R,M}$ denotes this mean when the first $J$
original core families are selected and the first $R$ protected. Along a joint
profile of a near-full sequence, $\meanarea_{\nu,J,R,M}$ is this mean for the
geometry of index $\nu$ on the subsequence of the profile, and
\begin{equation}\label{eq:gap-area-limit}
 \limarea_{R,M}=\limsup_{J\to\infty}\,\limsup_{\nu\to\infty}\,\meanarea_{\nu,J,R,M},
 \qquad \limarea_R=\limarea_{R,2R+2}.
\end{equation}
For $R<M$ both upper limits are real, and $0\le\limarea_{R,M}\le1/2$.
\end{definition}

\begin{theorem}[Theorem G: necessary frozen gap area]\label{thm:necessary-gap-area}
\formal{GapAreaMemo/GapArea.lean}{SourceCovers.gap_area_paid_deficit,SourceCovers.AdmissibleGeometry.prefix_gap_area_deficit,SourceCovers.NearFullSequence.JointProfile.prefix_gap_area_limit,SourceCovers.NearFullSequence.JointProfile.necessary_gap_area,SourceCovers.NearFullSequence.JointProfile.necessary_gap_area_default}
\begin{enumerate}[label=(\alph*)]
\item \emph{Finite form.} Take the data of \cite[Theorem~3.5]{SourceCovers2026}: a
fixed complete geometry with a nonempty correctness set, whole-step capacity at most
$\beta$ and an $N$-class span partition of its full family system; a $\mathrm Q_2$
selected core of at most $J$ families; at most $q$ exceptional families; $M\ge2$ and a
protected set of at most $R<M$ selected families; a bound $a$ on the hop count of
every family outside the selected core and the exceptions; and a real $\tau$ at least
the hop mass of the unprotected original nonleaves of the selected core. Then, for
every integer $w\ge1$,
\begin{equation}\label{eq:finite-gap-area}
 \frac\Delta L\ge\meanarea-\mathsf E_q(w;J,M,R,a,\tau).
\end{equation}
No runtime bound is used.
\item \emph{Prefixes.} For an admissible geometry with any order $G_1,G_2,\dots$ of its
original core by nonincreasing hop count, and integers $J\ge R\ge0$,
$M\ge\max(2,R+1)$, $w\ge1$ and $a\ge a_0$ with $a\ge cL/(J+1)$,
\[
 \frac\Delta L\ge\meanarea_{J,R,M}
 -\mathsf E_q\Bigl(w;J,M,R,a,\sum_{R<i\le J}|G_i|\Bigr).
\]
\item \emph{Dimension limit.} Along every joint profile of every near-full sequence,
for all integers $J\ge R\ge0$ and $M\ge\max(2,R+1)$ and every real $\lambda>0$,
\[
 \limsup_{\nu\to\infty}\meanarea_{\nu,J,R,M}
 \le\frac{\beta\tau_{R,J}}{M-R}+2\lambda+\frac{\beta NC_qc^2}{\lambda(J+1)}.
\]
\item \emph{The limit bound.} For every joint profile of every near-full sequence and
all integers $R\ge0$ and $M\ge\max(2,R+1)$,
\begin{equation}\label{eq:GA}
 \limarea_{R,M}\le\frac{\beta\tau_R}{M-R}.\tag{GA}
\end{equation}
In particular $\limarea_R\le\beta\tau_R/(R+2)$.
\end{enumerate}
In (c) and (d) the cutoff $R$ and the modulus $M$ are fixed, and the number $q$ of
exceptional families is any fixed integer.
\end{theorem}

\begin{outline}
(a) Apply \lean{SourceCovers.paid_deficit_of_gap_credit}, the formal paid inequality
of \cite[Theorem~3.5]{SourceCovers2026}, with the frozen point area
$\gaparea(S_\rho(s_j))$ as the credit of each legal phase $\rho$ and column $j$. Its
premise, that each credit is at most the frozen point area, holds with equality. The
mean credit is then $\meanarea$, and the paid error is unchanged.

(b) The first $J$ original families form a $\mathrm Q_2$ selection of at most $J$
families, and the hop mass of its unprotected nonleaves is at most
$\sum_{R<i\le J}|G_i|$. An omitted core family has at most $cL/(J+1)$ hops, by the
sorted order and the total bound $\sum_i|G_i|\le cL$; every other non-exceptional
family has at most $a_0$ hops. Hence (a) applies with this $a$, as in the proof of
\cite[Theorem~4.2]{SourceCovers2026}.

(c) For fixed $J,R,M,\lambda$, the error-limit lemma of the source-cover paper,
\lean{SourceCovers.NearFullSequence.JointProfile.prefix_error_limit}, supplies along
the subsequence widths $w_\nu\ge1$ and values $a_\nu$ admissible in (b), with widths
$\lceil\lambda L_\nu\rceil$ as in the proof of \cite[Theorem~4.2]{SourceCovers2026},
and upper bounds $F_\nu$ for the paid errors such that $\Delta_\nu/L_\nu+F_\nu$ tends
to the right side of~(c). By (b), $\meanarea_{\nu,J,R,M}\le\Delta_\nu/L_\nu+F_\nu$, and
the means lie in $[0,1/2]$; pass to the upper limit.

(d) Let $J\to\infty$ in (c): $\tau_{R,J}\to\tau_R$ by summability, and the last term
tends to zero. The finite upper limits lie in $[0,1/2]$, so the outer upper limit
preserves the inequality. Finally take $\lambda=\delta/2$ for every $\delta>0$ and
remove the additive $\delta$. For $M=2R+2$, $M-R=R+2$. As in
\cite[Theorem~4.6]{SourceCovers2026}, the dimension limit, the prefix limit and the
removal of the width error occur in that order.
\end{outline}

\paragraph{Why it matters.}
Theorem~G is the necessary condition that the paid inequality gives without choosing a
cover witness. For an impossibility proof it reduces the exclusion of a class of
admissible geometries to one lower bound: it suffices to show that every near-full
sequence in the class has a joint profile with $\limarea_{R,M}>\beta\tau_R/(M-R)$ for
some $R$ and $M$. Such a bound concerns the geometries alone, since frozen point areas
and hop profiles do not involve the correctness sets. Theorems~A and~B prove such
bounds for two classes. For a construction it is a requirement on the frozen banks.
For example, $\gaparea(S)\ge1/(2|S|)$ for every nonempty $S$ by the Cauchy--Schwarz
inequality. Along a joint profile of a near-full sequence with an inverse-square
limiting profile $h_i\le H/i^2$, for which $\tau_R\le H/R$ when $R\ge1$, the reciprocal
sizes $1/|S_\rho(s_j)|$ of the frozen raw sets therefore have limiting mean at most
$2\beta H/(R(R+2))$ at modulus $2R+2$ for every $R\ge1$, in the sense of
\eqref{eq:gap-area-limit}. This consequence is elementary and
is not formalized.

\begin{remark}[Comparison with the reciprocal-cover bound]\label{rem:reciprocal-comparison}
Proof sketch only: for every joint profile of every near-full sequence, and more
generally of every sequence of admissible geometries with $d\to\infty$ (the argument
does not use near-fullness, and $Q^*_{R,M}$ is defined along the profile as in
\cite[(4.7)]{SourceCovers2026}), and all integers $K,R\ge0$ and $M\ge\max(2,R+1)$,
\[
 \frac{(1-V_K)_+^2}{2(2K+1)}\,Q^*_{R,M}\le\limarea_{R,M},
\]
where $Q^*_{R,M}$ is the reciprocal-cover limit of \cite[(4.7)]{SourceCovers2026}.
Hence \eqref{eq:GA} implies the weighted reciprocal-cover bound
\cite[(4.8)]{SourceCovers2026}, and through it the necessary hop-tail bound
\cite[Theorem~4.2]{SourceCovers2026}. Indeed, for $J\ge\max(K,R)$ distinguish the first $K$
selected families. At every legal phase and column, \cite[Lemma~3.3]{SourceCovers2026}
gives
\[
 \gaparea(S_\rho(s_j))\ge\frac{(1-\alpha_\nu)_+^2}{2(2K+1)\kappa_{\nu,J,R,M,\rho,j}},
 \qquad \alpha_\nu=\sum_{K<i\le J}v_i^{(\nu)},
\]
with the fractional cover $\kappa$ of \cite[Definition~3.1]{SourceCovers2026}; its
formal form, \lean{SourceCovers.fractional_phase_gap}, is stated for the same point
area. Average over legal phases and columns. Along the profile
$\alpha_\nu\to\sum_{K<i\le J}v_i\le V_K$, and all averages lie in $[0,1]$, so the
dimension upper limit of $\meanarea_{\nu,J,R,M}$ is at least $(1-V_K)_+^2/[2(2K+1)]$
times that of the mean reciprocal cover. The prefix upper limit preserves this.
Theorem~G retains the full frozen point area, including what the fractional-cover
estimate gives up through the factors $(1-\alpha)_+^2$ and $2(2K+1)$.
\end{remark}
